
\documentclass{article}
\usepackage{colortbl}
\usepackage{makecell}
\usepackage{multirow}
\usepackage{supertabular}

\begin{document}

\newcounter{utterance}

\centering \large Interaction Transcript for game `eqbench', experiment `eq\_bench', episode 60 with reflexion{-}Qwen3.5{-}9B.
\vspace{24pt}

{ \footnotesize  \setcounter{utterance}{1}
\setlength{\tabcolsep}{0pt}
\begin{supertabular}{c@{$\;$}|p{.15\linewidth}@{}p{.15\linewidth}p{.15\linewidth}p{.15\linewidth}p{.15\linewidth}p{.15\linewidth}}
    \# & \multicolumn{2}{c}{Player} && \multicolumn{2}{c}{Game Master} \\
    \hline

    \theutterance \stepcounter{utterance}  
    & & & \multicolumn{4}{p{0.6\linewidth}}{
        \cellcolor[rgb]{0.9,0.9,0.9}{
            \makecell[{{p{\linewidth}}}]{
                \texttt{\tiny{[P1$\langle$GM]}}
                \texttt{You have attempted this same task before and did not succeed. These are the plans you wrote after those attempts.} \\
\texttt{Trial \#1: New plan: I incorrectly prioritized "Surprise" and "Shame" because Claudia was reacting to Tyler's accusations, but I failed to recognize that the core conflict is Claudia's defensive reaction to being labeled a liar and accused of abandoning Tyler. Tyler's aggressive tone ("lies," "abandoned me") directly triggers defensiveness and frustration, making "Anger" the primary emotion rather than a secondary one. I should have analyzed the dialogue's emotional trajectory: Claudia starts defensive, escalates when accused of lying, and ends in a state of justified outrage at his refusal to listen. Therefore, I will re{-}evaluate the scores to reflect that "Anger" is the dominant response to the confrontation, likely scoring it higher than "Shame" or "Surprise," while keeping "Anticipation" low as she is focused on the immediate conflict, not future events. Specifically, I will assign Anger a score of 9 or 10, Shame a score of 4 or 5 (as she feels guilty but is currently fighting back), and Surprise a score of 2 (she is not shocked by his anger, as she expected it).} \\
\\ 
\texttt{Trial \#2: New plan: I failed to correctly identify the primary emotion because I underestimated the intensity of Claudia's defensive rage in response to Tyler's specific accusations of lying and abandonment. While she feels shame for her past actions, the immediate context of being called a liar and having her motives dismissed as greed ("throwing money") triggers a stronger, more active emotional response: Anger. My previous scoring incorrectly balanced these, giving Shame a higher score than Anger. To fix this, I must recognize that "Save it" and "I'm not lying" are direct confrontations, not passive admissions. Therefore, I will re{-}evaluate the scores to reflect Anger as the dominant emotion (9{-}10), Shame as a significant but secondary underlying feeling (4{-}5), and keep Surprise and Anticipation very low (1{-}2) since the dialogue is a predictable, high{-}conflict confrontation rather than a revelation of new information or future planning. Specifically, I will assign Anger: 9, Shame: 4, Surprise: 1, Anticipation: 1.} \\
\\ 
\texttt{Trial \#3: New plan: I failed to recognize that Claudia's primary emotional state is defensive outrage rather than passive guilt. While she feels shame for her past actions, Tyler's aggressive accusations ("lies," "abandoned me," "throwing money") and his refusal to listen have triggered a fight response. Her repeated denials ("I'm not lying," "Save it") indicate she is actively fighting back against his narrative, which manifests as Anger. I should have prioritized Anger as the dominant emotion (9{-}10) because it drives her current behavior, while treating Shame as a secondary, underlying feeling (4{-}5) that is being suppressed by her defensiveness. Surprise and Anticipation remain low (1{-}2) as the confrontation is direct and focused on the present conflict, not future events or unexpected revelations. Therefore, I will assign Anger: 9, Shame: 4, Surprise: 1, Anticipation: 1.} \\
\\ 
\\ 
\\ 
\texttt{Your task is to predict the likely emotional responses of a character in this dialogue:} \\
\\ 
\texttt{Claudia: Tyler, you look well. Teaching suits you.} \\
\texttt{Tyler: Yeah, well, someone's gotta do the dirty work.} \\
\texttt{Claudia: I didn't mean it like that. I...} \\
\texttt{Tyler: I know what you meant. You always do, don't you?} \\
\texttt{Claudia: Not always, no. But I try to understand, at least.} \\
\texttt{Tyler: Understand? Like you understood when you left me alone?} \\
\texttt{Claudia: That's not fair, Ty. I had to...} \\
\texttt{Tyler: What? Chase your dreams? Leave your past behind?} \\
\texttt{Claudia: I never left you behind, Tyler. You were always...} \\
\texttt{Tyler: Always what? An afterthought? A burden?} \\
\texttt{Claudia: No, never a burden. You were my responsibility.} \\
\texttt{Tyler: Responsibility? Is that why you're here now? After all these years?} \\
\texttt{Claudia: I'm here because I...I wanted to see you.} \\
\texttt{Tyler: Or maybe you're here because you need something. Like always.} \\
\texttt{Claudia: That's not true, Tyler. I...} \\
\texttt{Tyler: Save it, Claudia. I've heard enough of your lies.} \\
\texttt{Claudia: I'm not lying, Tyler. I...} \\
\texttt{Tyler: You abandoned me, Claudia. You can't just waltz back into my life and pretend everything is okay.} \\
\texttt{Claudia: I'm not pretending, Tyler. I'm trying to make things right.} \\
\texttt{Tyler: Make things right? How? By throwing money at the problem?} \\
\texttt{Claudia: What are you talking about?} \\
\texttt{Tyler: Don't play dumb, Claudia. I know about the inheritance.} \\
\texttt{Claudia: I...I didn't know you knew.} \\
\texttt{Tyler: Of course you didn't. You never bothered to find out, did you?} \\
\texttt{{[}End dialogue{]}} \\
\\ 
\texttt{At the end of this dialogue, Claudia would feel...} \\
\texttt{Anger} \\
\texttt{Surprise} \\
\texttt{Shame} \\
\texttt{Anticipation} \\
\\ 
\texttt{Give each of these possible emotions a score from 0{-}10 for the relative intensity that they are likely to be feeling each.} \\
\\ 
\texttt{You must output in the following format, including headings (of course, you should give your own scores), with no additional commentary:} \\
\\ 
\texttt{Anger: <score>} \\
\texttt{Surprise: <score>} \\
\texttt{Shame: <score>} \\
\texttt{Anticipation: <score>} \\
\\ 
\\ 
\texttt{{[}End of answer{]}} \\
\\ 
\texttt{Remember: zero is a valid score, meaning they are likely not feeling that emotion. You must score at least one emotion > 0.} \\
\\ 
\texttt{Your answer:} \\
            }
        }
    }
     \\ \\

    \theutterance \stepcounter{utterance}  
    & \multicolumn{4}{p{0.6\linewidth}}{
        \cellcolor[rgb]{0.561,0.737,0.561}{
            \makecell[{{p{\linewidth}}}]{
                \texttt{\tiny{[P1$\rangle$GM]}}
                \texttt{Anger: 9} \\
\texttt{Surprise: 2} \\
\texttt{Shame: 4} \\
\texttt{Anticipation: 1} \\
            }
        }
    }
    & & \\ \\

    \theutterance \stepcounter{utterance}  
    & & & \multicolumn{2}{p{0.3\linewidth}}{
        \cellcolor[rgb]{0.9,0.9,0.9}{
            \makecell[{{p{\linewidth}}}]{
                \texttt{\tiny{[GM$|$GM]}}
                \texttt{Anger: 9} \\
\texttt{Surprise: 2} \\
\texttt{Shame: 4} \\
\texttt{Anticipation: 1} \\
            }
        }
    }
    & & \\ \\

    \theutterance \stepcounter{utterance}  
    & & & \multicolumn{2}{p{0.3\linewidth}}{
        \cellcolor[rgb]{0.9,0.9,0.9}{
            \makecell[{{p{\linewidth}}}]{
                \texttt{\tiny{[GM$|$GM]}}
                \texttt{Anger: 3} \\
\texttt{Surprise: 6} \\
\texttt{Shame: 8} \\
\texttt{Anticipation: 0} \\
            }
        }
    }
    & & \\ \\

    \theutterance \stepcounter{utterance}  
    & & & \multicolumn{2}{p{0.3\linewidth}}{
        \cellcolor[rgb]{0.9,0.9,0.9}{
            \makecell[{{p{\linewidth}}}]{
                \texttt{\tiny{[GM$|$GM]}}
                \texttt{game\_result = LOSE} \\
            }
        }
    }
    & & \\ \\

    \theutterance \stepcounter{utterance}  
    & & \multicolumn{4}{p{0.6\linewidth}}{
        \cellcolor[rgb]{0.95,0.95,0.95}{
            \makecell[{{p{\linewidth}}}]{
                \texttt{\tiny{[P1$\leftrightarrow$P]}}
                \texttt{Post{-}episode Reflexion (not sent to GM):} \\
\texttt{New plan: I have consistently failed to identify the correct primary emotion because I am over{-}indexing on Claudia's underlying guilt (Shame) and underestimating the immediate, reactive nature of her response to Tyler's aggression. The dialogue shows Claudia being systematically dismantled by Tyler's accusations ("lies," "abandoned me," "throwing money"). Her responses ("Save it," "I'm not pretending") are not passive admissions of guilt but active, defensive rebuttals. This indicates that her dominant state is not the quiet feeling of Shame, but the active, heated response of Anger triggered by being disbelieved and attacked. To succeed, I must re{-}evaluate the scores to reflect Anger as the overwhelming primary emotion (10), acknowledging that her defensiveness is a fight response, while reducing Shame to a secondary, suppressed feeling (3{-}4) that is currently being overridden by her outrage. Surprise and Anticipation should remain negligible (0{-}1) as the confrontation is direct and focused entirely on the present conflict. Specifically, I will assign Anger: 10, Shame: 3, Surprise: 0, Anticipation: 0.} \\
            }
        }
    }
    & \\ \\

\end{supertabular}
}

\end{document}
